\documentclass[10pt, a4paper]{article}
\usepackage[a4paper, top=2.2cm, bottom=2.2cm, left=2cm, right=2cm]{geometry}
\usepackage{fontspec}
\usepackage[english, polish, bidi=basic, provide=*]{babel}
\babelprovide[import, onchar=ids fonts]{polish}
\babelfont{rm}{Noto Sans}
\usepackage{array, tabularx, booktabs, colortbl, xcolor, fancyhdr, lastpage}

% ---------- Paleta (skala szarości) ----------
\definecolor{gDark}{gray}{0.22}
\definecolor{gMid}{gray}{0.55}
\definecolor{gHead}{gray}{0.80}
\definecolor{gLight}{gray}{0.92}

\setlength{\parindent}{0pt}
\setlength{\emergencystretch}{3em}
\renewcommand{\arraystretch}{1.35}
\arrayrulecolor{gMid}\setlength{\arrayrulewidth}{0.5pt}

% ---------- Nagłówek i stopka ----------
\pagestyle{fancy}
\fancyhf{}
\renewcommand{\headrule}{\hbox to\headwidth{\color{gMid}\leaders\hrule height 0.4pt\hfill}}
\renewcommand{\footrule}{\hbox to\headwidth{\color{gMid}\leaders\hrule height 0.4pt\hfill}}
\fancyhead[L]{\footnotesize\color{gDark} Bank Komercyjny S.A.}
\fancyhead[R]{\footnotesize\color{gDark} Dokument przekazywany do Organu Nadzorczego}
\fancyfoot[L]{\footnotesize\color{gDark} Kod referencyjny incydentu: DORA-INC-2026-004921-01}
\fancyfoot[R]{\footnotesize\color{gDark} Strona \thepage{} z \pageref{LastPage}}

% ---------- Makra formularza ----------
\newcommand{\sekcja}[1]{\par\vspace{9pt}\noindent\colorbox{gDark}{\parbox{\dimexpr\textwidth-2\fboxsep\relax}{\color{white}\bfseries\small #1}}\par\nopagebreak\vspace{1pt}}
\newenvironment{pola}
  {\noindent\tabularx{\textwidth}{|>{\columncolor{gLight}[\dimexpr\tabcolsep-0.4pt\relax][\dimexpr\tabcolsep-0.4pt\relax]\raggedright\arraybackslash\bfseries\small}p{6.2cm}|>{\raggedright\arraybackslash}X|}\hline}
  {\endtabularx\par}
\newcommand{\pole}[2]{#1 & #2 \\ \hline}
\newcommand{\podpis}[4]{\begin{minipage}[t]{0.47\textwidth}\raggedright\small\textbf{#1}\\[28pt]\rule{\linewidth}{0.4pt}\\[2pt]#2\par #3\par {\footnotesize #4}\end{minipage}}

\begin{document}

\begin{center}
    {\Large\bfseries WSTĘPNE POWIADOMIENIE O POWAŻNYM INCYDENCIE ZWIĄZANYM Z ICT}\\[4pt]
    {\small Przekazywane do Organu Nadzorczego (KNF) na podstawie art. 19 ust. 4 lit. a Rozporządzenia (UE) 2022/2554 (DORA)}\\[1pt]
    {\footnotesize Treść i terminy: Rozporządzenie delegowane (UE) 2025/301; wzory formularzy: Rozporządzenie wykonawcze (UE) 2025/302 (Załącznik I)}
\end{center}
\vspace{-4pt}\noindent\rule{\textwidth}{0.8pt}

% ============================================================
\sekcja{Ogólne informacje o podmiocie finansowym}
\begin{pola}
\pole{1.1 Rodzaj zgłoszenia}{Wstępne powiadomienie}
\pole{1.2 Nazwa podmiotu przedkładającego sprawozdanie}{Bank Komercyjny S.A.}
\pole{1.3 Kod identyfikacyjny podmiotu przedkładającego sprawozdanie}{004921}
\pole{1.4 Rodzaj podmiotu finansowego, którego dotyczy incydent}{Instytucja kredytowa}
\pole{1.5 Nazwa podmiotu finansowego, którego dotyczy incydent}{Bank Komercyjny S.A.}
\pole{1.6 Kod LEI podmiotu finansowego, którego dotyczy incydent}{2594000000BANKKOM049}
\pole{1.7 Imię i nazwisko głównej osoby wyznaczonej do kontaktów}{Marek Nowak}
\pole{1.8 Adres e-mail głównej osoby wyznaczonej do kontaktów}{soc@bank-komercyjny.example}
\pole{1.9 Numer telefonu głównej osoby wyznaczonej do kontaktów}{+48 22 000 11 22}
\pole{1.10 Imię i nazwisko dodatkowej osoby wyznaczonej do kontaktów}{}
\pole{1.11 Adres e-mail dodatkowej osoby wyznaczonej do kontaktów}{}
\pole{1.12 Numer telefonu dodatkowej osoby wyznaczonej do kontaktów}{}
\pole{1.13 Nazwa jednostki dominującej najwyższego szczebla}{}
\pole{1.14 Kod LEI jednostki dominującej najwyższego szczebla}{}
\pole{1.15 Waluta sprawozdawcza}{EUR}
\end{pola}

% ============================================================
\sekcja{Zawartość wstępnego powiadomienia}
\begin{pola}
\pole{2.1 Kod referencyjny incydentu nadany przez podmiot finansowy}{DORA-INC-2026-004921-01}
\pole{2.2 Data i godzina wykrycia poważnego incydentu związanego z ICT}{2026-11-09 14:30 UTC}
\pole{2.3 Data i godzina sklasyfikowania incydentu związanego z ICT jako poważny}{2026-11-09 15:15 UTC}
\pole{2.4 Opis poważnego incydentu związanego z ICT}{Awaria środowiska produkcyjnego usługi autoryzacji transakcji kartowych (niedostępność usługi wspierającej funkcję krytyczną lub ważną). Skutek operacyjny: niemożność realizacji płatności bezgotówkowych w punktach handlowo-usługowych oraz transakcji internetowych (e-commerce).}
\pole{2.5 Kryteria klasyfikacji, które doprowadziły do zgłoszenia incydentu}{Usługi krytyczne (art. 6); Klienci, kontrahenci finansowi i transakcje (art. 9 ust. 1); Wpływ reputacyjny (art. 9 ust. 2); Czas trwania i przestój usługi (art. 9 ust. 3); Wpływ ekonomiczny (art. 9 ust. 6)}
\pole{2.6 Progi istotności dla kryterium klasyfikacji „Zasięg geograficzny”}{Niespełnione -- wpływ ograniczony wyłącznie do terytorium Rzeczypospolitej Polskiej (PL); brak wpływu transgranicznego}
\pole{2.7 Wykrycie poważnego incydentu związanego z ICT}{Monitoring techniczny (alert zespołu SOC)}
\pole{2.8 Wskazanie, czy poważny incydent związany z ICT wywodzi się od dostawcy będącego osobą trzecią lub innego podmiotu finansowego}{Nie stwierdzono na obecnym etapie; dostawca oprogramowania systemowego został zawiadomiony}
\pole{2.9 Uruchomienie planu ciągłości działania, jeżeli plan został uruchomiony}{Tak -- uruchomiono procedurę Ciągłości Działania (BCP): przełączenie ruchu autoryzacyjnego na węzeł zapasowy}
\pole{2.10 Inne istotne informacje}{Brak przesłanek wskazujących na naruszenie poufności lub integralności danych. Data i czas przekazania powiadomienia: 2026-11-09 17:30 UTC (zachowano termin 4 h od klasyfikacji oraz 24 h od wykrycia).}
\end{pola}



% ============================================================
\sekcja{Osoby odpowiedzialne}
\begin{minipage}{\textwidth}
\vspace{4pt}
\podpis{Punkt Kontaktowy ICT}{Marek Nowak (SOC Team Lead)}{Tel.: +48 22 000 11 22}{E-mail: soc@bank-komercyjny.example}
\hfill
\podpis{Osoba Upoważniona}{Jan Kowalski (CISO)}{Dokument podpisany kwalifikowanym podpisem elektronicznym}{}
\end{minipage}

\end{document}